%% ============================================================
%%  Greylock Edge Application — Final Answers
%%  GlobalSushrut, Inc. — 2026
%% ============================================================
\documentclass[10pt,a4paper]{article}

\usepackage[T1]{fontenc}
\usepackage[utf8]{inputenc}
\usepackage{lmodern}
\usepackage{microtype}
\usepackage[top=0.75in,bottom=0.75in,left=0.85in,right=0.85in]{geometry}
\usepackage[dvipsnames,table]{xcolor}
\usepackage{enumitem}
\usepackage{setspace}
\usepackage[most]{tcolorbox}
\usepackage{titlesec}
\usepackage{fancyhdr}
\usepackage[
  colorlinks=false,
  pdftitle={Greylock Edge -- Application Answers},
  pdfauthor={GlobalSushrut Inc.},
]{hyperref}

%% ── Palette ──────────────────────────────────────────────────
\definecolor{CNavy}{RGB}{15,23,42}
\definecolor{CBlue}{RGB}{24,78,160}
\definecolor{CGray}{RGB}{100,116,139}
\definecolor{CSoft}{RGB}{246,248,252}
\definecolor{CLine}{RGB}{203,213,225}

%% ── Layout ───────────────────────────────────────────────────
\setstretch{1.08}
\setlength{\parindent}{0pt}
\setlength{\parskip}{5pt}
\setlist[itemize]{itemsep=2pt,topsep=2pt,leftmargin=1.4em}
\pagestyle{fancy}\fancyhf{}
\fancyhead[L]{\small\color{CGray}\textsc{Greylock Edge — Application Answers}}
\fancyhead[R]{\small\color{CGray}GlobalSushrut, Inc.\ \textbullet\ 2026}
\fancyfoot[C]{\small\color{CGray}\thepage}
\renewcommand{\headrulewidth}{0.35pt}

%% ── Question block ───────────────────────────────────────────
\tcbset{qbase/.style={
  enhanced, breakable, arc=4pt, boxrule=0pt,
  left=10pt, right=10pt, top=8pt, bottom=8pt,
  borderline west={3pt}{0pt}{CBlue}
}}

\newcommand{\qblock}[2]{%
  \vspace{8pt}
  \noindent
  \begin{tcolorbox}[qbase,
    colback=CSoft,
    title={\small\bfseries\color{CBlue} #1},
    attach boxed title to top left={yshift=-2mm,xshift=6pt},
    boxed title style={colback=CBlue,arc=2pt,boxrule=0pt,
      fonttitle=\small\bfseries\color{white}}]
  #2
  \end{tcolorbox}
}

%% ── Section style ────────────────────────────────────────────
\titleformat{\section}[runin]
  {\normalfont\large\bfseries\color{CBlue}}{}{}{}[\ \textcolor{CLine}{\rule[0.5ex]{0.6\linewidth}{0.4pt}}]
\titlespacing*{\section}{0pt}{14pt}{6pt}

%% ============================================================
\begin{document}
\pagestyle{fancy}

%% ── Mini header ──────────────────────────────────────────────
\begin{center}
  {\fontsize{8}{10}\selectfont\scshape\color{CGray} Greylock Edge Program Application}\\[3pt]
  {\fontsize{18}{22}\selectfont\bfseries\color{CNavy} Final Answers}\\[2pt]
  {\small\color{CGray}Copy-paste ready \textbullet\ One answer per question \textbullet\ No fluff}
\end{center}

\noindent\textcolor{CLine}{\rule{\linewidth}{0.6pt}}
\vspace{2pt}

%% ============================================================
%% Q3
%% ============================================================
\qblock{Q3 — Products you have built and launched, and your role}{

I built \textbf{Connector OS} — an AI execution control plane — from scratch over the past 18
months, as sole architect and primary engineer. Every layer is mine: the compiler, the runtime
enforcement engine, the cryptographic proof plane, the CLI, the REST API, and the operator UI.

The core of Connector is a \textbf{compiled governance layer} for AI agents.
I designed and implemented:

\begin{itemize}
  \item \textbf{CCL (Connector Contract Language)} — a domain-specific language and 7-pass
        compiler (lex → parse → semantic analysis → IR lowering → optimization → verification → emit)
        that compiles agent behavior into Ed25519-signed, content-addressed execution contracts.
        Policy violations are caught at compile time, not at runtime.

  \item \textbf{9-Chain Cryptographic Audit Architecture} — nine independent tamper-evident
        chains (audit, memory, dehallucination, compliance, governance, execution, trust, proof,
        isolation) that together produce a signed proof bundle covering HIPAA, SOC~2, GDPR,
        and EU AI Act from a single API call. Verifiable by a third party without a live system.

  \item \textbf{Namespace Isolation Model} — typed data classification (PHI in \texttt{/p/},
        knowledge in \texttt{/k/}, memory in \texttt{/m/}) enforced at the storage layer.
        Structural compliance — not a policy document.

  \item \textbf{Real-Time Budget Enforcement Gateway} — intercepts every LLM call, records
        provider-verified token counts per agent, enforces compile-time-declared budget ceilings.
        When the ceiling is hit, the agent stops with a signed receipt. No crash, no silence.
\end{itemize}

I also built live adversarial demos: prompt-injection attacks that Connector blocks mid-execution,
budget drain loops that self-terminate, and PHI leakage scenarios where the proof bundle shows
the LLM never received the sensitive data.

\textbf{This is not a project. It is infrastructure — a new category.}

}

%% ============================================================
%% Q4
%% ============================================================
\qblock{Q4 — Domains / sectors}{

\textbf{Infrastructure} \textbullet\ \textbf{Security}

\medskip
I build foundational layers, not applications.
Connector OS is to AI agents what a kernel is to processes — the mandatory
substrate through which all resource access, policy enforcement, and audit recording flows.
Selecting anything other than Infrastructure and Security would misrepresent what I am building.

}

%% ============================================================
%% Q5
%% ============================================================
\qblock{Q5 — How you landed on this problem space}{

I didn't choose this problem. I found it by building — and I couldn't un-see it.

I was building multi-agent systems when I noticed a structural absence that should not exist.
Every compute system ever built — operating systems, databases, cloud platforms — has a
\textbf{resource governance layer}. Processes don't access memory directly; they go through a kernel.
Database queries go through an access control layer. Network packets go through a firewall.

AI agents have none of this. They are the only class of software that can spend money, read sensitive
customer data, make consequential decisions, and produce external effects —
\textbf{with zero structural control, zero runtime enforcement, and zero cryptographic proof}.

The existing responses are wrong-layer answers:
LLM safety proxies filter prompts and outputs — that is like securing a server by checking the keyboard.
Orchestration frameworks (LangChain, CrewAI) add composition without adding governance.
Cloud AI safety tools sit at the model I/O boundary, not the execution layer.

The fundamental flaw I discovered:
\textit{we built sophisticated AI agency on top of an architecturally ungoverned substrate.}
The fix is not a feature. It is a new infrastructure layer — one that sits between AI agents and
everything they touch, enforcing compiled policy, recording cryptographic evidence,
and producing compliance proof that a regulator can verify without trusting the vendor.

I spent 18 months building that layer instead of writing about it.

}

%% ============================================================
%% Q6
%% ============================================================
\qblock{Q6 — Do you have a co-founder already?}{

\textbf{No} — and by design.

The first thing I needed to prove was that the architecture was correct.
You cannot co-found a system you have not fully thought through.
I built the core alone: the compiler, the proof plane, the runtime, the compliance layer.
I now know exactly what I built, what it takes to extend it, and what kind of partner would
compound it — not duplicate it.

I am looking for a co-founder whose strength is the enterprise GTM and relationship motion
in regulated industries. The technical architecture is decided. The moat is built.
The right partner runs toward healthcare, fintech, and legal AI buyers while I extend the platform.

}

%% ============================================================
%% Q7
%% ============================================================
\qblock{Q7 — Anything else you'd like us to know}{

I built the execution and evidence layer for AI systems —
the missing kernel that makes AI decisions controllable, auditable, and provable in production.

\medskip
This matters because regulated industries \textbf{cannot} deploy AI without it —
not by preference, but by law.
HIPAA, the EU AI Act, SOC~2 Type~II, and emerging SEC disclosure rules all require
exactly what Connector produces.
The window for voluntary adoption is closing. The mandate wave is arriving.

\medskip
Why me: I am one of very few people who has built a production AI governance system that is both
cryptographically correct and regulation-aware — not as a paper design, but as working, tested
infrastructure with a compiler, a proof chain, a live runtime, and adversarial demos.
\textbf{This is not a pitch. This is a system.}

}

%% ============================================================
%% Q8
%% ============================================================
\qblock{Q8 — Interest in Greylock talent team}{

\textbf{Yes} — selectively.

I am building a company. But Greylock portfolio companies operating at the intersection of
AI governance, security, and enterprise compliance are exactly the environments that sharpen
my understanding of the buyer. I am open to aligned conversations —
not because I am looking for a role, but because the right signal from the right environment
is how category companies get built faster.

}

\vspace{8pt}
\noindent\textcolor{CLine}{\rule{\linewidth}{0.5pt}}\\[3pt]
{\footnotesize\color{CGray}\centering
  Connector OS \textbullet\ GlobalSushrut, Inc. \textbullet\
  \textit{``The system that governs non-deterministic AI must itself be deterministic.''}
}

\end{document}
